% Article généré automatiquement
\documentclass[11pt,a4paper]{article}
\usepackage[utf8]{inputenc}
\usepackage[T1]{fontenc}
\usepackage[french]{babel}
\usepackage{graphicx}
\usepackage{amsmath}
\usepackage{hyperref}
\title{Reconnaissance faciale par LBP et fusion de similarités}
\author{Implémentation Java (OpenCV/Bytedeco, ImageJ)}
\date{\today}
\begin{document}
\maketitle
\begin{abstract}
Cet article décrit une chaîne complète de reconnaissance faciale implémentée en Java basée sur l'opérateur Local Binary Patterns (LBP), la construction d'histogrammes par blocs locaux, la fusion de plusieurs mesures de similarité et une décision binaire optimisée à partir de la courbe ROC.
\end{abstract}

\section{Introduction}
La reconnaissance faciale automatique repose sur plusieurs étapes clefs : localisation du visage, alignement ou recadrage, extraction de descripteurs discriminants et comparaison entre descripteurs. Notre implémentation s'appuie sur des méthodes classiques robustes pour des scénarios à contraintes modérées : détection par cascade, description texturale par LBP, et combinaison de mesures de similarité.

\section{Méthodes}
\subsection{Détection et recadrage}
Détection par classifieurs en cascade OpenCV. Les ROIs sont gérées en libérant explicitement les ressources natives.

\subsection{Prétraitement}
Conversion en niveaux de gris, égalisation locale et redimensionnement (W,H) pour grille stable.

\subsection{Extraction LBP}
Opérateur LBP, partition en grille MxN de blocs, histogrammes par bloc et concaténation. Voir annexe pour formules.

\subsection{Normalisation}
Normalisation L1 par bloc suivie d'une normalisation L2 globale pour certaines comparaisons.

\subsection{Mesures et fusion}
Chi-deux, similarité cosinus (convertie en distance), distance euclidienne. Fusion linéaire: $S_{fusion}=w_1\chi^2 + w_2 S_{cos} + w_3 d_2$.

\section{Protocole expérimental}
Benchmark séquentiel leave-one-out; extraction des caractéristiques sur disque puis comparaisons en flux.

\section{Résultats}
AUC approximative : 0.7064. Seuil choisi (max F1) : 61.0.

\begin{figure}[ht]
\centering
\includegraphics[width=0.7\textwidth]{target/figures/roc_curve.png}
\caption{Courbe ROC (AUC) obtenue.}
\end{figure}

\begin{figure}[ht]
\centering
\includegraphics[width=0.7\textwidth]{target/figures/score_distribution.png}
\caption{Distribution des scores (genuine vs impostor).}
\end{figure}

\section{Discussion}
Voir le fichier texte pour discussion, améliorations possibles et références.

\section{Conclusion}
Chaîne reproductible et recommandations pour amélioration (alignement, descripteurs profonds, calibration des scores).

\section*{Références}
\begin{itemize}
\item T. Ojala et al., Pattern Recognition, 1996.
\item P. Viola, M. Jones, CVPR, 2001.
\item M. Turk, A. Pentland, Journal of Cognitive Neuroscience, 1991.
\end{itemize}

\end{document}
